\begin{lemma}
\label{editorial-source-lemma-surjective-locally-nilpotent-kernel}
Let $\varphi : R \to S$ be a surjective map with locally nilpotent kernel.
Then $\varphi$ induces a homeomorphism of spectra and isomorphisms
on residue fields. For any ring map $R \to R'$ the ring map
$R' \to R' \otimes_R S$ is surjective with locally nilpotent kernel.
\end{lemma}

\begin{proof}
Put $I=\ker\varphi$ and retain the isomorphism
$R/I\to S$, $[r]\mapsto\varphi(r)$, supplied
by surjectivity. Every prime of $R$ contains
$I$, since each element of $I$ is nilpotent.
The inverse maps on primes are therefore
$$
\mathfrak q\longmapsto\varphi^{-1}(\mathfrak q),
\qquad \mathfrak p\longmapsto\varphi(\mathfrak p).
$$
The second is a prime ideal because its
quotient is the original domain $R/\mathfrak p$.
Surjectivity and $I\subseteq\mathfrak p$
show both composites are identities.
Moreover $D(\varphi(r))$ corresponds to
$D(r)$ for every original $r\in R$;
these sets form bases, so the bijection
and its inverse are continuous.

For corresponding primes $\mathfrak p$
and $\mathfrak q$, the residue-field map is
the explicit isomorphism
$$
\operatorname{Frac}(R/\mathfrak p)
\longrightarrow\operatorname{Frac}(S/\mathfrak q),
\qquad
\frac{r+\mathfrak p}{t+\mathfrak p}
\longmapsto
\frac{\varphi(r)+\mathfrak q}{\varphi(t)+\mathfrak q},
\quad t\notin\mathfrak p.
$$
The inverse lifts numerator and denominator
along $\varphi$; different lifts differ by
$I\subseteq\mathfrak p$, so the fractions
are unchanged. A nonzero denominator on
either side remains nonzero on the other.

For an arbitrary original map $f:R\to R'$,
there are inverse algebra maps
$$
R'\otimes_R S\longrightarrow R'/IR',\qquad
a\otimes\varphi(r)\longmapsto[af(r)],
\qquad
R'/IR'\longrightarrow R'\otimes_R S,\quad
[a]\longmapsto a\otimes1.
$$
Changing a lift $r$ of an element of $S$
adds an element of $I$, and balancing
uses the same $f(r)$, so the first map
is well-defined. The second kills $IR'$
because $f(i)\otimes1=1\otimes\varphi(i)=0$
for $i\in I$. The formulas prove both
composites are identities. Thus the
original map $R'\to R'\otimes_R S$ is
surjective with kernel exactly $IR'$.

Each element of that kernel is a finite
sum $\sum_{j=1}^m a_j f(i_j)$ with $i_j\in I$.
Choose $d_j\geq1$ with $i_j^{d_j}=0$.
In the full multinomial expansion of its
$D$th power, for $D=1+\sum_j(d_j-1)$,
every term contains some $f(i_j)^{d_j}$,
and is zero. Thus $IR'$ is locally
nilpotent, without a uniform exponent
being assumed for the whole ideal.
Applying the first part to this actual
quotient proves the asserted base-change
conclusions, with the same prime and
residue-field maps.

The kernel hypothesis is also necessary
for a surjective ring map to induce a
homeomorphism onto the entire spectrum:
its image is $V(I)$, so surjectivity on
primes requires $I$ to lie in every
prime, equivalently in the nilradical.
In addition, the given quotient induces
an exact isomorphism of reduced rings
$R_{red}\to S_{red}$. Indeed if
$\varphi(r)$ is nilpotent, then $r^n\in I$
for some $n$, and $(r^n)^m=0$ for some
$m$, so $r$ is nilpotent. The converse
is immediate under a ring homomorphism.
The preimage of $\operatorname{Nil}(S)$
is consequently exactly
$\operatorname{Nil}(R)$, proving the
stated reduced-quotient isomorphism and
the same assertion after each base change.
\end{proof}